\documentclass[12pt]{article}
\usepackage{amsmath,amssymb,amsthm}
\usepackage[margin=1in]{geometry}
\usepackage{algorithm}
\usepackage{algpseudocode}
\usepackage{booktabs}

\newtheorem{theorem}{Theorem}
\newtheorem{lemma}[theorem]{Lemma}
\theoremstyle{definition}
\newtheorem{definition}[theorem]{Definition}

\title{Problem 161: Triominoes}
\author{Project Euler Solution}
\date{}

\begin{document}
\maketitle

\section{Problem Statement}
A triomino is a shape consisting of three unit squares joined edge-to-edge. There are two combinatorial types: the I-triomino (three collinear squares, 2 orientations) and the L-triomino (three squares in an ``L,'' 4 orientations), giving 6 placement types in total. Determine the number of distinct tilings of a $9 \times 12$ grid by triominoes.

\section{Definitions}

\begin{definition}
A \emph{tiling} of a rectangular grid $R$ by triominoes is a partition of the unit cells of $R$ into triominoes, where each cell belongs to exactly one piece.
\end{definition}

\begin{definition}
The \emph{broken profile} at cell $(r,c)$ during a row-major traversal is the bitmask encoding which cells ahead of $(r,c)$ in the traversal order are already occupied by previously placed triominoes.
\end{definition}

\section{Mathematical Development}

\begin{theorem}[Existence and finiteness]
The number of distinct tilings of the $9 \times 12$ grid by triominoes is a well-defined non-negative integer.
\end{theorem}

\begin{proof}
The grid has $9 \times 12 = 108$ cells. Each triomino covers exactly 3 cells, and $3 \mid 108$, so the necessary divisibility condition is satisfied. The set of valid tilings is a finite subset of all partitions of 108 cells into groups of 3.
\end{proof}

\begin{lemma}[Placement geometry]
Each of the 6 triomino placement types, anchored at the topmost-leftmost cell in row-major order, occupies cells at the following offsets $(dr, dc)$:
\begin{center}
\begin{tabular}{@{}lccc@{}}
\toprule
Type & Offsets & Row span & Column span \\
\midrule
L-A  & $(0,0),\,(0,1),\,(1,0)$ & 2 & 2 \\
L-B  & $(0,0),\,(0,1),\,(1,1)$ & 2 & 2 \\
L-C  & $(0,0),\,(1,0),\,(1,1)$ & 2 & 2 \\
L-D  & $(0,0),\,(1,-1),\,(1,0)$ & 2 & 2 \\
I-H  & $(0,0),\,(0,1),\,(0,2)$ & 1 & 3 \\
I-V  & $(0,0),\,(1,0),\,(2,0)$ & 3 & 1 \\
\bottomrule
\end{tabular}
\end{center}
\end{lemma}

\begin{proof}
By direct enumeration. Each L-triomino is a $2 \times 2$ block minus one corner.
\end{proof}

\begin{theorem}[Broken-profile DP correctness]
Processing the grid cell-by-cell in row-major order, the state at each cell is fully characterized by a bitmask of width at most $W = 12$ bits (the number of columns). The tiling count equals $\mathrm{dp}[0]$ after processing the final cell.
\end{theorem}

\begin{proof}
At cell $(r,c)$, any anchored triomino occupies cells whose row-major offset from $(r,c)$ is at most $2 \cdot 12 = 24$. After shifting out the current cell at each step, the profile requires at most $\mathrm{COLS}$ bits. If the current cell is filled (bit~0 set), we skip; otherwise we try all 6 placements whose cells lie in-bounds with unoccupied bits, marking them. After all 108 cells, $\mathrm{dp}[0]$ counts exactly the valid tilings.
\end{proof}

\begin{lemma}[State space bound]
The number of reachable profiles at any step is at most $2^{12} = 4096$, but empirically much smaller due to consistency constraints.
\end{lemma}

\section{Algorithm}

\begin{algorithmic}[1]
\Function{CountTilings}{$\mathrm{ROWS}=9, \mathrm{COLS}=12$}
    \State $\mathrm{dp}[0] \gets 1$
    \For{$r = 0$ \textbf{to} $\mathrm{ROWS}-1$}
        \For{$c = 0$ \textbf{to} $\mathrm{COLS}-1$}
            \State $\mathrm{ndp} \gets \text{empty map}$
            \ForAll{$(\mathrm{mask}, w)$ in dp}
                \If{$\mathrm{mask} \mathbin{\&} 1 = 1$}
                    \State $\mathrm{ndp}[\mathrm{mask} \gg 1] \mathrel{+}= w$
                \Else
                    \ForAll{placement $P$ in 6 types}
                        \If{$P$ fits in bounds and all bits free}
                            \State $\mathrm{ndp}[\mathrm{nmask} \gg 1] \mathrel{+}= w$
                        \EndIf
                    \EndFor
                \EndIf
            \EndFor
            \State $\mathrm{dp} \gets \mathrm{ndp}$
        \EndFor
    \EndFor
    \State \Return $\mathrm{dp}[0]$
\EndFunction
\end{algorithmic}

\section{Complexity Analysis}
\begin{itemize}
    \item \textbf{Time:} $O(\mathrm{ROWS} \cdot \mathrm{COLS} \cdot |\mathcal{S}| \cdot 6)$ where $|\mathcal{S}| \le 2^{12} = 4096$. Worst case: $O(108 \cdot 4096 \cdot 6) \approx O(2.7 \times 10^6)$.
    \item \textbf{Space:} $O(|\mathcal{S}|) = O(2^{12})$ for the DP hash map at each step.
\end{itemize}

\section{Answer}

\[
\boxed{20574308184277971}
\]

\end{document}
